\chapter{SYSTEM IMPLEMENTATION}

\section{Introduction}
The implementation phase of the project involves the translation of the design specifications into a functional software and hardware system. This phase ensures that the system components—frontend, backend, database, and hardware—interact seamlessly to fulfill the project requirements.

\section{Software Environment}
The software development environment was carefully selected to ensure high performance and scalability.
\begin{itemize}
    \item \textbf{Frontend}: React.js was used for building the dynamic user interface. 
    \item \textbf{Backend}: Node.js with Express.js framework handles the business logic and API requests.
    \item \textbf{Database}: MongoDB Atlas (Cloud) provides the NoSQL storage for staff and attendance records.
    \item \textbf{Firmware}: Arduino IDE with ESP32 board support for BLE scanning.
\end{itemize}

\section{Frontend Development}
The frontend is built using a component-based architecture. Key components include:
\subsection{Dashboard Component}
The dashboard provides a real-time summary of staff presence. It uses React hooks (useEffect and useState) to fetch and update data from the backend every 5 seconds.
\subsection{Attendance Reports}
This module allows administrators to export attendance data into various formats. It includes filtering options by department, date range, and staff ID.

\section{Backend Logic and API Routes}
The backend acts as the bridge between the BLE detections and the user interface.
\subsection{BLE Synchronizer}
This route (\texttt{/api/ble-data}) receives data from multiple ESP32 nodes. It performs deduplication to ensure that if a staff member is detected by two receivers, only the strongest signal is recorded.

\section{Hardware Integration}
The hardware setup involves placing ESP32 receivers at a height of 2 meters to avoid signal absorption by humans. 
\begin{itemize}
    \item \textbf{Detection Cycle}: Scanning (5s) -> Averaging (1s) -> Transmitting (1s).
    \item \textbf{Filtering}: All signals with RSSI weaker than -95dBm are ignored to prevent detection from outside the room.
    \item \textbf{Power Management}: The ESP32 is powered via the classroom's existing 5V USB adapters. It remains always on, drawing less than 0.5W of power.
    \item \textbf{Dual Mode Capability}: The hardware operates concurrently as a BLE Scanner and a GATT Server. This means it listens for passive tags while accepting active connections from mobile phones if physical tags are forgotten.
\end{itemize}

\section{Security Implementation}
Security is heavily enforced at both the hardware and software layers. 
\begin{itemize}
    \item API routes require JWT authentication using the Bearer token scheme. 
    \item Role-based access control prevents regular staff from accessing the administrative substitution routes.
    \item Hardware spoofing is mitigated by matching the UUIDs to the encrypted database records.
\end{itemize}
